%!TeX program = xelatex
\documentclass[12pt,hyperref,a4paper,UTF8,fontset=fandol]{ctexart}
\usepackage{SDUReport}
\usepackage{booktabs}
\usepackage{amsmath}
\usepackage{float}
\usepackage{graphicx}
\usepackage{hyperref}

\begin{document}

\cover
\thispagestyle{empty}

\newpage
\pagenumbering{Roman}
\tableofcontents

\newpage
\pagenumbering{arabic}

\section{任务背景与要求}

本项目完成两段工业流水线视频中的装配件缺失检测。视频前段只包含正常装配样本，测试段同时包含正常件与异常件。系统需要在不使用异常样本训练的条件下，对每个稳定到位的工件输出 \texttt{OK} 或 \texttt{NG}，并在演示视频中显示 ROI、检测分数与状态。

作业指南要求包含三部分产物：预处理、训练和推理代码；能够展示缺盖案例的短演示视频；以及说明方法、结果和改进思路的 PDF 报告。本项目的代码入口为 \texttt{tools/run\_demo.py}，核心模块位于 \texttt{src/}，最终演示视频输出到 \texttt{results/}。

\section{总体方案}

系统采用无监督异常检测路线。训练阶段只从正常装配片段中提取稳定工件的 ROI patch，使用预训练 DINOv2 ViT-S/14 提取 patch token 特征，并用 PatchCore 记忆库进行最近邻异常评分。测试阶段按工件分组，对稳定区间采样评分，并聚合为单件 OK/NG 结果。

整体流程如下：

\begin{enumerate}
    \item 跳过视频前 60 秒抖动片段。
    \item 根据固定 ROI 或动态定位得到检测区域。
    \item 通过运动、模糊、前景与存在性门控，只在工件稳定到位时检测。
    \item 对训练段正常 ROI 提取 DINOv2 patch 特征，建立 PatchCore memory bank。
    \item 对测试段每件工件计算异常分数，超过阈值判定为 \texttt{NG}。
    \item 使用仪表盘视频显示当前状态、分数、ROI 框与 OK/NG 结果。
\end{enumerate}

\section{预处理与分件}

\subsection{稳定工件检测}

流水线视频中存在空场景、进入和离开过程、运动模糊以及工件只露出一部分的情况。若逐帧检测，会产生大量误报。因此本项目实现了按件检测：当工件进入 ROI 后，\texttt{UnitTracker} 根据帧差、前景比例、模糊度和最小稳定帧数建立一个工件单元；只有单元稳定时才送入检测器。

对于无工件或半进入状态，系统只显示 \texttt{WAIT} 或 \texttt{MOVING}，不输出 OK/NG，从而符合实际质检中“到位后判定”的要求。

\subsection{Task A 动态定位}

Task A 中圆形端面位置会有轻微漂移。项目在固定参考 ROI 的基础上加入动态定位模块 \texttt{src/locator.py}，在每帧中定位端面区域并裁剪对齐 patch。该处理降低了位置漂移对 PatchCore 分数的影响。

Task A 当前使用动态 ROI、DINOv2 特征和 PatchCore 判定，并额外通过端面颜色比例辅助识别盖子缺失后的黄色内衬区域。

\subsection{Task B 多帽位 ROI}

Task B 的一个工件包含 4 个关键帽位。项目为 4 个帽位分别设置 ROI，并在视频旋转 180 度后同步旋转 ROI。4 个帽位角色为：蓝、白、白、蓝。系统要求至少 3 个帽位可见且稳定时才进入检测，避免半进入或半离开工件被误判。

\section{核心算法}

\subsection{DINOv2 特征提取}

项目使用 \texttt{timm} 加载 \texttt{vit\_small\_patch14\_dinov2}。每个 ROI 被缩放到统一输入尺寸，送入模型后取 patch token 作为局部特征。与手工颜色阈值相比，DINOv2 特征能同时表达边缘、结构、纹理和部件形态；与有监督检测器相比，它不需要异常标注。

\subsection{PatchCore 记忆库}

训练段中的正常 patch 特征构成候选特征集合。为降低存储和 kNN 搜索开销，系统使用 coreset 子采样压缩 memory bank。测试时，对每个 patch 特征计算其到正常 memory bank 最近邻的距离，并按 \texttt{mean\_top5} 聚合为 ROI 异常分数。

阈值由训练段正常分数自适应估计：

\begin{equation}
    \theta = \mu_{train} + k \cdot \sigma_{train}
\end{equation}

其中 $k$ 由配置文件中的 \texttt{threshold\_k} 控制。

\subsection{Task B 独立 ROI PatchCore}

Task B 若把四个帽位混在同一个 memory bank 中，会削弱每个帽位的局部差异。因此最终方案对四个 ROI 分别训练 PatchCore，每个帽位拥有独立阈值。单件检测时，只要任一 ROI 的聚合分数超过对应阈值，就判定该工件为 \texttt{NG}。

此外，Task B 针对蓝色和白色帽位分别计算颜色存在性比例。颜色存在性不作为主判定，而作为 PatchCore 的兜底信号：只有持续、明显低于训练正常分布时才触发 NG。最终配置中 \texttt{taskB\_presence\_ng\_vote\_ratio=0.90}，避免轻微反光或白色区域波动造成误报。

\section{实验结果}

\subsection{演示视频}

最终演示视频位于：

\begin{itemize}
    \item \texttt{results/taskA\_demo\_result.mp4}
    \item \texttt{results/taskB\_demo\_result.mp4}
\end{itemize}

演示视频包含工件状态、OK/NG 结果、异常分数、阈值和 ROI 框。空场景与运动过程显示为等待或移动状态，不作为质检结果计数。

\subsection{Task A 结果}

Task A 使用动态定位后的单 ROI PatchCore。训练段仅使用正常装配样本，测试时按件聚合。经过动态定位后，端面漂移导致的误报显著减少；盖子缺失时，端面纹理和颜色变化会使 PatchCore 分数升高，并被判定为 \texttt{NG}。

\subsection{Task B 结果}

Task B 最终使用 \texttt{1040s--1075s} 附近片段作为短演示窗口，该窗口包含两个明显异常工件。50 秒源视频回归时，系统输出 10 件，其中 8 件 \texttt{OK}、2 件 \texttt{NG}；两个 \texttt{NG} 分别对应 1047 秒和 1072 秒附近的异常工件。生成 15 秒提交演示时，使用同一异常区间的 2 倍速短片。

\begin{table}[H]
    \centering
    \begin{tabular}{lccc}
        \toprule
        片段 & 工件数 & OK & NG \\
        \midrule
        Task B 1040s--1090s 回归片段 & 10 & 8 & 2 \\
        Task B 250s--300s 回归片段 & 8 & 6 & 2 \\
        Task B 930s--980s 压力测试片段 & 11 & 8 & 3 \\
        \bottomrule
    \end{tabular}
    \caption{Task B 回归测试结果}
\end{table}

\section{速度与工程实现}

系统按 \texttt{step\_test=3} 抽帧检测，并对一个工件内部的稳定帧进行采样聚合，避免逐帧重复计算。PatchCore 使用 coreset 压缩 memory bank，减少 kNN 搜索成本。演示视频以 2 倍速输出，当前配置下 Task B 30 秒源视频生成约 15 秒演示视频，满足短视频展示要求。

主要代码结构如下：

\begin{itemize}
    \item \texttt{src/preprocessing.py}：视频读取、旋转、ROI 裁剪、训练样本提取。
    \item \texttt{src/tracker.py}：工件稳定状态机与按件分组。
    \item \texttt{src/locator.py}：Task A 动态端面定位。
    \item \texttt{src/feature\_extractor.py}：DINOv2 特征提取。
    \item \texttt{src/memory\_bank.py}：PatchCore memory bank、coreset 和 kNN 评分。
    \item \texttt{src/detection.py}：单 ROI 与多 ROI 异常检测器。
    \item \texttt{src/visualization.py}：演示视频仪表盘与 ROI 可视化。
    \item \texttt{tools/run\_demo.py}：训练、测试、可视化视频生成入口。
\end{itemize}

\section{困难情况与改进}

本任务中最难处理的是运动过程、反光、半进入工件和多帽位差异。Task A 的主要问题是端面位置漂移，因此加入动态定位。Task B 的主要问题是四个帽位外观不同，因此采用 per-ROI PatchCore，而不是把所有帽位混成一个模型。

后续可继续改进的方向包括：

\begin{itemize}
    \item 将 DINOv2 特征缓存到磁盘，减少重复生成演示视频时的训练时间。
    \item 使用 FAISS 或 ONNX/TensorRT 加速最近邻搜索和特征推理。
    \item 为 Task B 增加少量可解释的结构模板，仅用于展示异常位置，不替代 PatchCore 主判定。
    \item 引入自动评估脚本，对已知异常时间点计算召回率和误报率。
\end{itemize}

\section{结论}

项目完成了从视频预处理、ROI 定位、无监督模型训练、按件推理到结果可视化的完整流程。Task A 通过动态定位提升了稳定性，Task B 通过多 ROI PatchCore 解决了四个帽位差异明显的问题。最终产物满足作业指南中对代码、演示视频和报告的要求，并保持训练阶段只使用正常样本。

\end{document}
